\chapter{Conclusion}
\setcounter{equation}{0}
Phishing URLs remain a common entry point for attacks that seek to steal information, defraud users, or distribute malware. Machine-learning-based URL screening can help identify suspicious links, but performance measured on one dataset may not transfer to other URL populations.

This seminar reviewed X-PHIDE, the XGBoost-based URL phishing detection method proposed by Misiek and Hyla. The method combines allowlist and blocklist checks with engineered URL features, cross-dataset feature selection, Bayesian hyperparameter optimisation, XGBoost classification, and browser-based risk display. It uses three datasets: a custom set of 75,738 URLs, GramBeddings with 639,723 URLs, and a 47,103-URL feature-compatible subset of PhiUSIIL.\cite{bib1}

The reported model achieved 90.72\% accuracy and a 91.16\% F1 score on the custom test set. However, F1 scores fell to 73.25\% on GramBeddings and 69.37\% on PhiUSIIL, with an average of 77.93\%. XGBoost trained substantially faster than the neural network in the paper's benchmark, while the neural network offered faster inference and a smaller model size. The importance of HTTPS also varied considerably between datasets.

The main lesson is that detection quality, computational efficiency, and generalisation must be considered together. Stable feature selection and independent testing are essential, but they do not eliminate all risks. A practical system should be monitored, updated as phishing patterns evolve, and assessed for false alarms and missed threats. X-PHIDE provides a useful approach to browser-based phishing detection, while its cross-dataset results clearly demonstrate the need for continued validation.\cite{bib1}
